
\section{Test 3: Síntesi: nombres naturals i enters}

\iniciTest{tC}{b,b,c,a,d,c,b,a,c,a}

\begin{enumerate}

\pregunta Quina és la manera canònica d'identificar un nombre natural \(n\) amb un enter?
\begin{enumerate}
\opcio{a}{Mitjançant la classe \(\classe{(0,n)}\).}
\opcio{b}{Mitjançant la classe \(\classe{(n,0)}\).}
\opcio{c}{Mitjançant la classe \(\classe{(n,n)}\).}
\opcio{d}{Mitjançant la classe \(\classe{(0,0)}\).}
\end{enumerate}
\feedback

\pregunta Quina de les equacions següents té solució en \(\Z\), però no en \(\N\)?
\begin{enumerate}
\opcio{a}{\(x+2=5\).}
\opcio{b}{\(x+5=2\).}
\opcio{c}{\(x+0=7\).}
\opcio{d}{\(x+3=3\).}
\end{enumerate}
\feedback

\pregunta Si \(a,b\in\N\) i \(a\leq b\), què podem dir de l'enter \(\classe{(a,b)}\)?
\begin{enumerate}
\opcio{a}{És sempre positiu.}
\opcio{b}{És sempre nul.}
\opcio{c}{És no positiu, és a dir, \(\classe{(a,b)}\leq 0\).}
\opcio{d}{No es pot comparar amb \(0\).}
\end{enumerate}
\feedback

\pregunta Quina igualtat expressa correctament la resta en \(\Z\)?
\begin{enumerate}
\opcio{a}{\(a-b=a+(-b)\).}
\opcio{b}{\(a-b=a+b\).}
\opcio{c}{\(a-b=b-a\).}
\opcio{d}{\(a-b=-(a+b)\).}
\end{enumerate}
\feedback

\pregunta Si \(p,q\in\N\setminus\{0\}\), quin és el resultat de
\[
\classe{(0,p)}\cdot\classe{(0,q)}?
\]
\begin{enumerate}
\opcio{a}{Un enter negatiu.}
\opcio{b}{La classe \(\classe{(0,pq)}\).}
\opcio{c}{La classe \(\classe{(p,q)}\).}
\opcio{d}{La classe \(\classe{(pq,0)}\), que representa un enter positiu.}
\end{enumerate}
\feedback

\pregunta Quina diferència hi ha entre els divisors de \(12\) en \(\N\) i en \(\Z\)?
\begin{enumerate}
\opcio{a}{En \(\Z\), \(12\) no té divisors.}
\opcio{b}{En \(\N\), \(12\) té divisors negatius.}
\opcio{c}{En \(\Z\), també apareixen els oposats dels divisors positius.}
\opcio{d}{En \(\N\) i en \(\Z\) els divisors són exactament els mateixos.}
\end{enumerate}
\feedback

\pregunta En el teorema de la divisió en \(\Z\), si \(a=bq+r\), amb \(b\neq0\), quina condició ha de complir el residu?
\begin{enumerate}
\opcio{a}{\(-|b|<r<|b|\).}
\opcio{b}{\(0\leq r<|b|\).}
\opcio{c}{\(0<r\leq |b|\).}
\opcio{d}{\(r<0\).}
\end{enumerate}
\feedback

\pregunta Quina afirmació és correcta sobre el màxim comú divisor en el context dels enters?
\begin{enumerate}
\opcio{a}{Es pren habitualment el representant positiu; per exemple, \(\mcd(18,24)=6\).}
\opcio{b}{Sempre es pren el representant negatiu; per exemple, \(\mcd(18,24)=-6\).}
\opcio{c}{No té sentit parlar de màxim comú divisor en \(\Z\).}
\opcio{d}{El màxim comú divisor de dos enters sempre és \(0\).}
\end{enumerate}
\feedback

\pregunta Si \(n\in\N\setminus\{0\}\), què representa l'enter \(\classe{(0,n)}\)?
\begin{enumerate}
\opcio{a}{El natural \(n\).}
\opcio{b}{L'enter \(0\).}
\opcio{c}{L'enter negatiu \(-n\).}
\opcio{d}{Un nombre que no pertany a \(\Z\).}
\end{enumerate}
\feedback

\pregunta Quin mètode és natural per demostrar que tot nombre natural \(n\geq2\) és primer o producte de primers?
\begin{enumerate}
\opcio{a}{La inducció forta.}
\opcio{b}{La definició de l'oposat en \(\Z\).}
\opcio{c}{La divisió euclidiana amb divisor negatiu.}
\opcio{d}{La definició de valor absolut.}
\end{enumerate}
\feedback

\end{enumerate}

